% ============================================================
% PROPOSAL: Euclidean spinor / OS conventions block
% Version: v1 (2026-06-10). STATUS: DRAFT FOR GPT REVIEW.
% The LAST blocker before the Q11 implementation package.
% Deliberately short: conventions only, no new results.
% Proposed placement: new \paragraph at the start of
% \S sec:qs_reflection_positivity (where the scalar reflection
% \Theta is already defined), with a one-sentence forward pointer
% from \S sec:dirac_eq.
% ============================================================

\paragraph{Euclidean spinor conventions and the graded OS form.}
For the spinorial extension of the OS programme the following
conventions are fixed once and used throughout.
\emph{(i) Gammas.}
Euclidean gamma matrices are hermitian and satisfy
\begin{equation}
  \{\gamma_A,\gamma_B\}=2\delta_{AB},
  \qquad
  \gamma_A^{\dagger}=\gamma_A,
  \qquad A,B\in\{0,1,2,3\},
  \label{eq:q11_eucl_gamma}
\end{equation}
with $\gamma_0$ associated to the distinguished direction $w$
(equivalently the ordered-branch time, $t=w$).
For the mostly-plus Lorentzian convention used in the ordered
branch, $\eta_{\mu\nu}=\mathrm{diag}(-1,1,1,1)$
(as fixed earlier in the article), a compatible Wick map to the
Lorentzian matrices of \S\ref{sec:dirac_eq} is
\begin{equation}
  \gamma^0_L = i\gamma_0,\qquad
  \gamma^i_L=\gamma_i,
  \label{eq:q11_wick_map}
\end{equation}
up to the harmless overall sign convention for $\gamma^0_L$.
With this choice
$\{\gamma^\mu_L,\gamma^\nu_L\}=2\eta^{\mu\nu}$.
This relation is a convention-matching statement; the Euclidean OS
calculation below uses only the hermitian Euclidean matrices
$\gamma_A$.
\emph{(ii) Kernel.}
The free minimal Euclidean Dirac operator and two-point kernel are
\begin{equation}
  D_E=\gamma_A\partial_A+m,\qquad m>0,
  \qquad
  S_E=D_E^{-1},
  \qquad
  \tilde S_E(p)=\frac{m-i\gamma_Ap_A}{p^2+m^2},
  \label{eq:q11_eucl_dirac_kernel}
\end{equation}
with the Berezin pairing
$\langle\psi_\alpha(x)\bar\psi_\beta(y)\rangle
=S_{E,\alpha\beta}(x-y)$ for Grassmann fields $\psi$, $\bar\psi$.
\emph{(iii) Reflection.}
The OS reflection is the same $\Theta:(w,\vec x)\mapsto(-w,\vec x)$
already used in the scalar subsectors; on spinor fields it acts
antilinearly, reverses operator order, and carries the spinorial
reflection matrix $\gamma_0$ for the Dirac variables used here
(if the spinor sector is later rewritten in Majorana or
charge-conjugated variables, the reflection matrix and adjoint must
be translated accordingly):
\begin{equation}
  \Theta\psi(x)=\bar\psi(\Theta x)\,\gamma_0,
  \qquad
  \Theta\bar\psi(x)=\gamma_0\,\psi(\Theta x),
  \qquad
  \Theta(AB)=\Theta(B)\,\Theta(A).
  \label{eq:q11_eucl_theta_spinor}
\end{equation}
The Euclidean adjoint of a smeared field polynomial is defined
through this $\Theta$-operation together with complex conjugation
of coefficients.
\emph{(iv) Graded OS form.}
For elements $F,G$ of the field algebra supported in the half-space
$w>0$, the OS sesquilinear form is
\begin{equation}
  (F,G)\;\equiv\;\bigl\langle F\cdot\Theta G\bigr\rangle_E ,
  \label{eq:q11_graded_os_form}
\end{equation}
with the operator order fixed as written; for Grassmann fields the
two possible orderings differ by the reordering sign, and
\eqref{eq:q11_graded_os_form} is the order convention used in the
spinorial OS test below.
Its positivity properties are not assumed here; they are checked
explicitly in the free-kernel calculation.
The scalar reflection-positivity formula of this section is the
restriction of \eqref{eq:q11_graded_os_form} to the bosonic subalgebra.
\emph{(v) Scope and status.}
These conventions lie in the standard Euclidean Fermi-field OS
family (Osterwalder--Schrader Fermi fields; transfer-matrix
constructions), with $\gamma_0$ here playing the role usually
denoted $\gamma_4$.
They are fixed at the level of the free minimal first-order kernel;
Majorana/charge-conjugated reformulations, higher-derivative
kernels, and the interacting sector would require separate
treatment.
The conventions themselves carry no derivational content
(definition-level; no Level assignment); all statements proved with
them inherit their own Levels.

% ------------------------------------------------------------
% Forward pointer to add in \S sec:dirac_eq (one sentence):
% ------------------------------------------------------------
% "The Euclidean counterpart of these conventions (hermitian
% gammas, the Euclidean Dirac kernel, and the graded OS reflection
% structure) is fixed in \S\ref{sec:qs_reflection_positivity} and
% used in the spin--statistics analysis."
%
% ------------------------------------------------------------
% Bookkeeping notes (for the implementation package):
% - new labels: eq:q11_eucl_gamma, eq:q11_eucl_dirac_kernel,
%   eq:q11_eucl_theta_spinor, eq:q11_graded_os_form
% - symbol list additions: gamma_A (Euclidean), D_E, S_E, Theta on
%   spinors, (F,G) graded OS form, sigma (reordering sign — only if
%   the appendix calculation is inserted)
% - no bibliography additions strictly required; optional: OS 1973
%   Fermi-fields reference if not already present
% ============================================================
